\section{Related Work}
\label{sec:related}

Constrained QAOA spans penalty encodings and algorithms that restrict the
state support or dynamics. The alternating-operator framework formalizes
problem-aware initial states and feasible-subspace mixers
\cite{hadfield2019alternating}; broader constructions cover constraint
encodings and preserving mixers \cite{fuchs2022mixers,ruan2023constraints}.
Grover mixers transfer part of the burden to feasible-state preparation
\cite{bartschi2020grover}, Zeno dynamics uses repeated projections
\cite{herman2023zeno}, and indicator-function methods evaluate inequalities in
ancillas \cite{bucher2025penaltyfree}. These are alternatives to, rather than
matched variants of, the uniform full-space Penalty-X design studied here.

Recent structured approaches sharpen that contrast. Onah, Firt, and
Michielsen introduce a constraint-enhanced one-hot construction with $W$
states and block-$XY$ mixing \cite{onahfirt2025encoded}; related analysis
studies shallow feasibility limits for permutation constraints
\cite{onah2025limitations}. Bucher et al. combine iterative warm starts with
one-hot $XY$ mixers \cite{bucher2026iterative}, extending the wider warm-start
and initial-state/mixer-alignment literature
\cite{egger2021warm,he2023alignment}. RCSP flow and resource feasibility is
not a simple fixed-Hamming-weight constraint, so those constructions are not
treated as automatic RCSP baselines.

Classical loss design provides a different intervention. CVaR emphasizes a
low-energy tail \cite{barkoutsos2020cvar}; validation-oracle objectives can
target feasibility directly \cite{li2026feasibilityloss}; and feasibility-
driven penalty scheduling jointly varies constraint penalties
\cite{ferrari2026scheduling}. Onah and Michielsen give finite-depth,
finite-shot guarantees for a Fej\'er-filtered constrained model with explicit
phase-normalization assumptions \cite{onah2026fejer}. These studies motivate
careful separation of objective information, spectral assumptions, and
sampling resources. Our O2 instead serves only as an exact-statevector
capacity control, while O3 retains the same $H_C$ and ansatz as O0.

Finally, QAOA initialization and continuation are standard responses to
outer-loop difficulty \cite{zhou2020qaoa}; general training hardness and
reachability results warn against identifying a terminal local-optimizer
result with the global variational optimum
\cite{bittel2021training,akshay2021reachability}. The exact zero-layer test is
narrower: it certifies only that a reported deeper point is worse than a known
embedded point. The contribution of this study is the combined attribution
design---scale control, nested optimizer diagnosis, O2 capacity measurement,
and graph-held-out O3 confirmation---not a new ansatz, CVaR definition, or
generic optimizer-hardness theorem. Extended study-by-study comparisons are
in Online Resource~1, Tables~S14--S15.
